\section{Detailed Hardware Design}
\label{sec:hardware-design}

\subsection{Alfred Central Compute (ACC)}
\label{sec:acc}

\subsubsection{Processor architecture}

The ACC is a heterogeneous compute module, not a single SoC choice, because "run the application layer" and "close a control loop deterministically" are different workloads that should not share a scheduler.

\begin{itemize}
\item \textbf{Application processor (AP):} automotive-grade multi-core Arm (Cortex-A) SoC class, e.g., NXP S32G (also has integrated deterministic switch + lockstep cores), Qualcomm Snapdragon Ride / Cockpit-class, or Texas Instruments TDA4/AM69 class for the prototype-to-production span. For prototyping specifically, an NXP S32G274A or S32G399A evaluation board is the pragmatic starting point because it already integrates a PFE (packet forwarding engine) and lockstep Cortex-M7 cores alongside the A53 cluster --- i.e., it is a single part that models the AP/RT split this document requires.
\item \textbf{Real-time processor (RT):} lockstep or dual-core-with-comparator Cortex-R52/M7-class MCU, either integrated on the same SoC (S32G's M7 cores) or a discrete companion MCU (e.g., Infineon AURIX TC3xx, or ST SPC58) for programs that want ISO~26262 ASIL-D qualified silicon independent of the AP's qualification story. The RT processor hosts: network time distribution, safety-relevant supervisory logic, and any control loop that must survive an AP crash/reboot (e.g., hazard lights, brake light illumination logic staying alive even if the Linux side is mid-OTA-reboot).
\item \textbf{Ethernet switch:} either the SoC-integrated switch (S32G PFE, or a Marvell/NXP automotive Ethernet switch such as 88Q5050-class or SJA1110-class parts) providing hardware-accelerated L2 switching, VLAN segmentation per zone, and TSN (IEEE 802.1Qbv/Qav) traffic shaping so that control traffic is not starved by telemetry/logging traffic on the same backbone.
\end{itemize}

\subsubsection{Safety partitioning and OS strategy}

\begin{diagram}
 +---------------------------------------------------------------+
 |                     ALFRED CENTRAL COMPUTE                    |
 |                                                                 |
 |  +----------------------+   +--------------------------------+ |
 |  |  Application Domain  |   |   Real-Time / Safety Domain     | |
 |  |  (Cortex-A cluster)  |   |   (lockstep Cortex-R/M7)        | |
 |  |----------------------|   |----------------------------------| |
 |  | Linux (or hypervisor |   | Alfred RT Runtime (bare-metal /  | |
 |  |  + Linux guest)      |   |  small RTOS: Zephyr/FreeRTOS)    | |
 |  |                      |   |                                  | |
 |  | - Alfred Runtime     |   | - Safety-relevant supervisory    | |
 |  | - Machine IR store   |<->| logic (watchdog arbitration,     | |
 |  | - Behavior/App layer |   |   fail-safe outputs)             | |
 |  | - Deployment Manager |   | - Deterministic network time     | |
 |  | - Diagnostics/OTA    |   |   distribution (gPTP grandmaster)| |
 |  | - Logging/telemetry  |   | - Hard-real-time control loops   | |
 |  +----------------------+   |   too fast/critical to run       | |
 |                              |   at the edge (rare; see 5.1.3)  | |
 |                              +--------------------------------+ |
 |             |                              |                    |
 |             +---------------+--------------+                    |
 |                             |                                    |
 |                   +---------+---------+                          |
 |                   |  Integrated TSN   |                          |
 |                   |  Ethernet Switch  |                          |
 |                   +----+----+----+----+                          |
 |                        |    |    |                                |
 +------------------------|----|----|--------------------------------+
        100/1000BASE-T1   |    |    |  10BASE-T1S trunk / zone spurs
                    Zone A     Zone B    Zone C ...
\end{diagram}

\textbf{OS strategy:} Linux (Yocto-based automotive-grade BSP) on the AP for the application domain --- this is where the Alfred Runtime, Machine IR store, deployment manager, and diagnostics live, because these need a rich filesystem, networking stack, package management, and container-style isolation between application "skills." A type-1 hypervisor (e.g., Xen-based automotive hypervisors, or leveraging the S32G's hardware partitioning) is introduced only when a program requires running an OEM-supplied guest OS/instrument-cluster stack alongside Alfred on the same silicon; it is not required for the core thesis and should not be built until a customer requires co-residency. The RT domain runs a small deterministic RTOS (Zephyr is the pragmatic default: automotive contributors, native support for time-sensitive-networking stacks, permissive license) or bare-metal, and is treated as a separate safety element with its own boot chain.

\subsubsection{The placement framework (central vs. edge)}

This is the decision procedure the firmware generator (Section~\ref{sec:firmware-generation}) executes automatically for every control loop in the Machine IR:

\begin{center}
\small
\begin{tabularx}{\linewidth}{@{}l X X@{}}
\toprule
\textbf{Factor} & \textbf{Push toward central} & \textbf{Push toward edge} \\
\midrule
Deadline & $>$ 10--20~ms (round-trip network + scheduling budget on typical 10/100BASE-T1 + switch hop) & $<$ 1~ms, or any deadline where a single dropped/delayed frame is unacceptable (e.g., commutation timing) \\
Jitter tolerance & Loose (user-facing state changes, mode changes) & Tight (PWM commutation, current-loop sampling) \\
Bandwidth & Low-rate commands / setpoints & High-rate raw sensor streams that would saturate the backbone if centralized (e.g., 20~kHz current samples) \\
Fault containment & Loop failure should be visible/loggable centrally & Loop failure must not depend on network availability (network cable unplugged should not spin a motor away) \\
Safety integrity & Below ASIL A / QM, or where central redundancy (ACC RT domain) already provides the integrity & ASIL B+ functions requiring a local, independently-verifiable control loop \\
Network latency \& availability & Backbone is available $>$99.99\% of the time in normal operation & Any function that must degrade gracefully to a safe state on network loss \\
\bottomrule
\end{tabularx}
\end{center}

Worked examples used throughout this document: \code{spoiler.set\_position(65\%)} is a low-rate setpoint with a $\sim$500~ms acceptable settle time --- the trajectory/ramp planning executes centrally, and only the innermost position servo loop (closing on the Hall/potentiometer feedback at $\sim$1~kHz) executes on the E2B-MOTOR node. A 20~kHz current controller for a BLDC never leaves the edge node: it runs on the MCU's dedicated timer/ADC/PWM hardware with no network involvement in the loop at all --- the network only carries the current-loop's \emph{setpoint} and periodic telemetry summaries (mean current, fault flags), not every sample.

\subsubsection{ACC subsystem inventory}

\begin{itemize}
\item \textbf{Storage:} eMMC or automotive-grade UFS for the OS/application image (A/B partitioned for OTA); separate small SPI-NOR for the RT domain's boot image and for a minimal recovery/failsafe image on the AP side, independent of the main OS partitions.
\item \textbf{Secure boot:} hardware root of trust (SoC-integrated HSM/CAAM on S32G, or discrete Secure Element such as an automotive-grade SE05x/OPTIGA) verifying a chain: boot ROM $\rightarrow$ first-stage bootloader $\rightarrow$ U-Boot $\rightarrow$ kernel+initramfs (dm-verity rootfs) $\rightarrow$ Alfred Runtime container images (signed). RT domain has its own independent secure boot chain rooted in the same hardware root of trust but verified separately, so a compromise of the AP OS cannot re-flash RT firmware.
\item \textbf{Watchdogs:} hardware watchdog in the RT domain supervises AP liveness (heartbeat over an internal bus, e.g., SPI or shared memory + interrupt) and forces defined fail-safe outputs (e.g., hazard lights on, all actuators to last-known-safe or de-energized state) if the AP stops heartbeating; a second watchdog inside the AP OS supervises the Alfred Runtime process tree itself and can restart individual containerized "skills" without a full reboot.
\item \textbf{Power management:} automotive PMIC with cold-crank/load-dump tolerant input stage (see Section~\ref{sec:bom}), always-on RT domain rail sized for the minimum fail-safe function set, switched/sleep-capable AP rail so the vehicle can enter a low-power state without losing RT supervisory function.
\item \textbf{Redundant power:} for higher-integrity programs (UGV, agricultural machinery operating unattended), dual input rails from separate vehicle power sources feeding an ORing power stage into the ACC, and a supercapacitor or small buffer battery sized to ride through transients and to guarantee an orderly shutdown/log-flush on primary power loss.
\item \textbf{Diagnostics:} UDS-compatible diagnostic server on the AP for compatibility with existing shop tools, plus an Alfred-native diagnostics/telemetry channel (structured, versioned, richer than UDS DIDs) for Alfred's own tooling.
\item \textbf{OTA:} A/B image update with cryptographic signature verification, staged rollout support, and automatic rollback on boot-health-check failure (see Section~\ref{sec:deployment-reprogramming}).
\item \textbf{Logging:} local ring-buffered structured logs (per-subsystem) persisted to storage, mirrored to the Alfred backend when connectivity exists; all log entries timestamped against the same synchronized clock the Machine IR uses for correlation (Section~\ref{sec:probe-discovery}).
\item \textbf{Time synchronization:} the ACC RT domain acts as gPTP (IEEE 802.1AS) grandmaster for the vehicle Ethernet network; all E2B nodes are gPTP-synchronized slaves. This is not optional --- it is what allows telemetry from independent E2B nodes to be correlated and what a firmware/discovery pipeline depends on for causality inference (shared with Option B's synchronization requirement, Section~\ref{sec:probe-discovery}).
\item \textbf{Service/debug interface:} dedicated debug connector (USB-C for host-mode debug console + Ethernet-over-USB for a bench link independent of the vehicle network) physically separate from any customer-facing connector, disabled/gated in production builds behind a signed-challenge unlock rather than left permanently open.
\end{itemize}

\subsection{Alfred E2B Edge Node Family}
\label{sec:e2b-family}

\subsubsection{Universal node vs. a family of nodes}

\begin{decision}[title={Recommendation: a small family, not one universal node}]
Build a family of 4--5 standardized boards sharing one MCU family, one Ethernet PHY/transceiver strategy, one firmware base image, and one connector philosophy, but differing in the analog/power front-end. A single "do everything" node forces every deployment to pay for motor-driver silicon, isolation, and connector pins it doesn't use, which is both a BOM and a size/thermal penalty in a harness environment where connector count and enclosure size matter as much as the die area of the MCU. The commonality lives in \emph{firmware and software}, not silicon: every variant runs the same base image, same discovery/manifest agent, same OTA client, same secure boot chain, and exposes the same style of capability manifest --- only the driver layer and populated I/O differ. This is exactly the HAL discipline described in Section~\ref{sec:software-design}, pushed into the hardware SKU strategy itself.
\end{decision}

\subsubsection{Common architecture across all E2B variants}

\begin{diagram}
        +-----------------------------------------------------+
        |                  ALFRED E2B NODE                     |
        |                                                       |
        |   +------------------+      +---------------------+  |
        |   |   MCU (Arm        |      |  Ethernet PHY/MAC   |  |
        |   |  Cortex-M7/M33,    |<---->|  10BASE-T1S (LAN8670|  |
        |   |  automotive grade) |  RMII/  or 100/1000BASE-T1 |  |
        |   |                    |  MII    (DP83TC81x/       |  |
        |   |  - MAC integrated  |      |   TJA1103-class))   |  |
        |   |    or via PHY link |      +----------+----------+  |
        |   +---+--+--+--+--+---+                 |             |
        |       |  |  |  |  |                       |             |
        |     PWM ADC SPI UART GPIO           T1S/T1 line        |
        |       |  |  |  |  |                  transformer/      |
        |    (variant-specific analog/driver   AEC-qualified      |
        |     front end -- see 5.2.2)          coupling network   |
        |                                                          |
        |   +------------------+   +------------------+           |
        |   | Local flash/RAM   |   | Watchdog + Secure|           |
        |   | (firmware, config,|   | Boot (HSM or MCU |           |
        |   |  cached manifest) |   |  crypto engine)   |           |
        |   +------------------+   +------------------+           |
        |                                                          |
        |   +------------------+   +------------------+           |
        |   | Power supply      |   | Protection        |           |
        |   | (12/24V in, buck  |   | (reverse polarity, |           |
        |   |  to 3.3V/5V rails)|   |  TVS, load dump,   |           |
        |   |                    |   |  per-channel fuse/ |           |
        |   |                    |   |  eFuse on outputs) |           |
        |   +------------------+   +------------------+           |
        +-----------------------------------------------------+
                     |                                  |
              Debug: SWD/JTAG + UART           Field connector:
              header, populated at             sealed automotive connector
              manufacturing test only          (signal + power combined or
              (not exposed in final                separate, per variant)
              enclosure)
\end{diagram}

MCU family recommendation for prototyping through early production: STMicroelectronics STM32 automotive-grade parts on the H7/G4 family for compute-heavy variants (motor control, sensor fusion) and the automotive-qualified STM32 "A" suffix parts where AEC-Q100 is required, or NXP S32K3-series (purpose-built for exactly this: automotive MCU with integrated CAN-FD, LIN, and a security/HSM block). The S32K3 family is the stronger long-term bet for E2B specifically because NXP already ships an automotive-qualified Ethernet-capable variant, integrated CAN-FD/LIN, and a secure boot/HSM (EdgeLock) block in the same part, reducing BOM count relative to STM32 + discrete CAN/LIN transceivers only.

\subsubsection{Variant-specific front ends}

\begin{center}
\small
\begin{tabularx}{\linewidth}{@{}p{2.1cm} p{2.3cm} X X@{}}
\toprule
\textbf{Variant} & \textbf{Network PHY} & \textbf{Analog/driver front end} & \textbf{Representative use} \\
\midrule
E2B-LITE & 10BASE-T1S & GPIO bank (8--16 ch, protected), PWM out (4--8 ch), ADC (8 ch, 12-bit), no motor driver silicon & Switches, indicator lamps, simple relays, door-ajar sensors, seatbelt switches \\
E2B-MOTOR & 100BASE-T1 & 3-phase or H-bridge gate driver + FETs (rated to target current, e.g., 20--40A), current-sense amplifiers (shunt + INA-class), quadrature/Hall decoder inputs (dedicated timer channels), overcurrent comparator with hardware shutdown path independent of MCU & Window/seat/mirror/spoiler/HVAC-damper motors, small BLDC pumps \\
E2B-SENSOR & 10BASE-T1S (100BASE-T1 option) & ADC front end (14--16 bit for precision analog), SPI/I2C master (multi-drop sensor bus), SENT decoder (hardware timer capture or dedicated SENT peripheral) & Pressure, temperature, humidity, position sensors, occupant sensing \\
E2B-BODY & 10BASE-T1S & LIN transceiver (1--2 ch, e.g., TJA1027/1028-class), CAN-FD transceiver (1--2 ch, e.g., TCAN1462/TJA1462-class) for supplier sub-assemblies, GPIO, PWM (servo-style) & Bridging a purchased CAN/LIN sub-assembly (seat module, mirror module) into the Alfred network; also the natural home for Option B's legacy-adapter gateway role (Section~\ref{sec:deployment-reprogramming}) \\
E2B-HIGH-SPEED & 1000BASE-T1 (100BASE-T1 fallback) & High-speed serial/parallel camera or sensor interface (MIPI-CSI to Ethernet bridge, or native Ethernet-out sensor), larger local buffer RAM/flash for burst buffering & Cameras, lidar/radar front-ends, any peripheral whose native data rate exceeds what 10BASE-T1S or a 100BASE-T1 spur can sustain alongside other zone traffic \\
\bottomrule
\end{tabularx}
\end{center}

Every variant shares: automotive temperature range components (AEC-Q100/Q200 grade for production; commercial/industrial grade tolerated in prototype-only builds, called out explicitly in the BOM, Section~\ref{sec:bom}), a hardware watchdog independent of the MCU's own windowed watchdog (belt-and-suspenders against a firmware bug hanging the watchdog-service task itself), reverse-polarity and load-dump protection on the power input, per-output overcurrent protection (resettable eFuse preferred over a blown fuse in the field), a debug header populated only at manufacturing test and not present on the field connector, and a manufacturing/provisioning interface (Section~\ref{sec:probe-discovery}/\ref{sec:deployment-reprogramming}) used once at first boot to inject a unique device identity and manufacturing test results into secure storage.

Connector strategy: a single sealed automotive connector family (e.g., TE MCON or similar IP67-rated automotive connector series) standardized across the whole E2B family for network+power (2-pair or 4-pair depending on PoDL current needs), with a second, variant-specific connector for the peripheral-facing I/O, so that the network/power harness segment is identical part-to-part across a vehicle regardless of which E2B variant sits at that location --- this is what makes an E2B node field-replaceable without harness rework, one of the concrete "hardware becomes interchangeable" claims in the thesis.

Power-over-line consideration: 10BASE-T1S nodes are strong candidates for Power over Data Line (PoDL, IEEE 802.3bu) to reduce the sealed-connector pin count further; this should be prototyped early (MVP1/2, Section~\ref{sec:roadmap}) because it changes the connector and power-protection design non-trivially and is much cheaper to decide once than to retrofit.

\subsection{Alfred Probe Hardware (Option B)}
\label{sec:probe-hardware}

\subsubsection{Why distributed probes beat a single OBD-II dongle}

\begin{itemize}
\item \textbf{Topology visibility:} OBD-II is downstream of a gateway; many body/comfort/chassis signals never reach the diagnostic connector unfiltered, and switched automotive Ethernet segments are frequently not bridged to OBD-II at all. A probe node physically spliced onto CAN-B, another onto CAN-C, and a TAP on the automotive Ethernet backbone see traffic OBD-II structurally cannot.
\item \textbf{Simultaneity for causality inference:} the correlation techniques in Section~\ref{sec:probe-discovery} (Part 2-F, physical correlation) require observing a command frame on one bus and a physical effect (current draw, Hall pulse) at another physical point \emph{at the same instant}. A single dongle cannot be in two electrical locations at once; a distributed, clock-synchronized probe set can.
\item \textbf{Bus-load and error-state visibility:} a single point of observation cannot distinguish "this signal doesn't exist on this bus" from "this signal exists on a different bus segment behind a gateway that filters it." Multiple simultaneous vantage points resolve that ambiguity directly instead of by inference.
\end{itemize}

\subsubsection{Modular probe node design}

\begin{diagram}
      Physical interfaces (interchangeable front-end modules)
      +---------+  +---------+  +---------+  +---------+
      | CAN/    |  | LIN     |  |Ethernet |  | Analog/ |
      | CAN-FD  |  | module  |  | TAP     |  | GPIO/   |
      | module  |  |         |  | module  |  | PWM/enc |
      +----+----+  +----+----+  +----+----+  +----+----+
           |             |            |            |
           +------+------+------+-----+------+-----+
                  |                          |
           +------+------+           (high-speed path
           |  Front-end   |            for 100/1000BASE-T1
           |  aggregation |            TAP -- may bypass
           |  (FPGA, see  |            FPGA and go direct
           |  5.3.2)      |            to SoC MAC, see text)
           +------+------+
                  |
        Timestamped event stream (shared internal bus/PCIe/
        high-speed SPI depending on channel count)
                  |
           +------+------+
           | ARM/Linux SoC|  (e.g., NXP i.MX8M, or automotive-
           | probe host   |   grade equivalent for field units)
           +------+------+
                  |
           Gigabit Ethernet uplink to Alfred Analyzer
           (isolated network, not the vehicle network)
\end{diagram}

Modules: CAN/CAN-FD (isolated transceiver front end, listen-only mode capable at the transceiver level, not just software), LIN (isolated transceiver, master/slave-passive listen), 10BASE-T1S tap, 100BASE-T1 tap, 1000BASE-T1 tap, conventional Ethernet tap (passive optical/electrical tap or a monitor-port-capable switch), UART/RS-485, digital input bank, PWM capture (dedicated timer-capture channels, not polled GPIO, to get accurate duty-cycle/frequency), analog input (current-limited, buffered), current measurement (isolated Hall-effect or shunt+isolated amplifier), and pulse/quadrature encoder capture.

\subsubsection{Is an FPGA worth it?}

\begin{center}
\small
\begin{tabularx}{\linewidth}{@{}X X@{}}
\toprule
\textbf{FPGA front end (small automotive-adjacent FPGA, e.g., Lattice ECP5/CrossLink or similar low-power part)} & \textbf{No FPGA (SoC peripherals / MCU co-processor only)} \\
\midrule
Pros: deterministic hardware timestamping of every edge/frame independent of OS jitter; can host many parallel CAN/LIN/PWM-capture channels without contending for a shared MCU bus; reconfigurable channel mix (same board, different bitstream, becomes an 8-channel CAN probe or a mixed CAN+LIN+PWM probe) & Pros: much simpler firmware/tooling (no HDL team needed), faster iteration for a small startup, lower BOM cost and power at low channel counts, easier to source and qualify \\
Cons: requires FPGA/HDL expertise the team must build or hire; adds a bitstream build/verification pipeline parallel to the firmware pipeline; higher BOM cost and power at low channel counts & Cons: OS/MCU-side timestamping jitter (typically tens of microseconds on a well-tuned RTOS interrupt, worse on Linux without care) may be insufficient for tight causality inference on fast signals; scaling channel count means scaling MCU count \\
\bottomrule
\end{tabularx}
\end{center}

\begin{decision}[title={FPGA recommendation}]
Do not put an FPGA in the MVP probe (Section~\ref{sec:roadmap}, MVP0/1). Use a small automotive-grade MCU (e.g., an STM32 or S32K with hardware CAN-FD/LIN peripherals and dedicated timer-capture units) per 4--8 channel module feeding timestamped frames to the host SoC over a fast link, with the MCU's own free-running timer cross-calibrated against the shared clock (Section~\ref{sec:probe-discovery}). Revisit an FPGA front end once the probe needs to scale past roughly 16--24 simultaneous mixed-protocol channels per unit, or once sub-microsecond cross-channel timestamp correlation becomes a proven customer requirement rather than a hypothetical one --- introducing FPGA/HDL engineering before that need is validated is the single highest-risk "build too much, too early" mistake available in this architecture.
\end{decision}

For the high-bandwidth Ethernet TAP path specifically (100BASE-T1/1000BASE-T1), route captured frames directly into the host SoC's own Ethernet MAC(s) via passive TAP hardware (a purpose-built automotive TAP, or two PHYs in monitor configuration) rather than through the FPGA/MCU aggregation path --- that path already has hardware timestamping support (PTP hardware timestamping is a standard feature of automotive Ethernet MACs) and forcing gigabit-class traffic through a low-pin-count aggregation bus would become the bottleneck.
